% !TEX program = pdflatex
% !TEX encoding = UTF-8 Unicode

\documentclass[12pt,letterpaper]{article}

% Paquetes esenciales
\usepackage[utf8]{inputenc}
\usepackage[T1]{fontenc}
\usepackage[spanish,es-tabla]{babel}
\usepackage{lmodern}
\usepackage{microtype}

% Geometría y diseño
\usepackage[left=2.5cm,right=2.5cm,top=2.5cm,bottom=2.5cm]{geometry}
\usepackage{setspace}
\onehalfspacing

% Tablas y gráficos
\usepackage{booktabs}
\usepackage{longtable}
\usepackage{array}
\usepackage{pdflscape}

% Hipervínculos y metadatos
\usepackage[hidelinks]{hyperref}

% Encabezados y pies de página
\usepackage{fancyhdr}
\pagestyle{fancy}
\fancyhf{}
\fancyhead[L]{\small Derecho Penal Especial Mexicano}
\fancyhead[R]{\small S8: Juego de roles}
\fancyfoot[C]{\thepage}
\renewcommand{\headrulewidth}{0.4pt}

% Definiciones de metadatos
\def\documentauthor{Mart\'in Jonathan de la Cruz Mu\~noz}
\def\documentsubject{Licenciatura en Derecho}
\def\documenttitle{Actividad S8: Juego de roles -- Delitos contra la paz, seguridad de las personas e inviolabilidad del domicilio}

\hypersetup{
    pdftitle={\documenttitle},
    pdfauthor={\documentauthor},
    pdfsubject={\documentsubject},
    pdfkeywords={Derecho Penal, Amenazas, Cobranza ilegítima, Allanamiento, Usurpación de identidad},
    pdfproducer={LaTeX con pdfTeX},
    pdfcreator={AulaTeX Academic System}
}

% Configuración del título
\title{\textbf{\documenttitle}}
\author{
    \documentauthor\\
    \small Matrícula: ES2611202040\\
    \small \documentsubject\\
    \small Universidad Abierta y a Distancia de México (UnADM)
}
\date{\small Semestre 2 -- Bloque 2\\[0.5em] \today}

\begin{document}

\maketitle
\thispagestyle{fancy}

\vspace{1em}

\section{Introducción}

Los delitos contra la paz, la seguridad de las personas y la inviolabilidad del domicilio constituyen un conjunto de figuras típicas que protegen bienes jurídicos fundamentales en el orden social y la convivencia pacífica. La legislación penal mexicana, tanto en el ámbito federal como en las entidades federativas, contempla diversos tipos penales destinados a salvaguardar la tranquilidad individual, la integridad patrimonial y moral de las personas, así como la privacidad del hogar.

El presente trabajo analiza cuatro figuras delictivas de especial relevancia: las amenazas, la cobranza ilegítima, el allanamiento de morada y la usurpación de identidad. Estas conductas representan diferentes modalidades de afectación a la seguridad y paz de los individuos, cada una con características particulares en su configuración típica y en los elementos probatorios requeridos para su acreditación.

El marco normativo aplicable se encuentra contenido principalmente en el Código Penal Federal (CPF) y en el Código Penal para el Distrito Federal (CPDF), hoy Ciudad de México. El estudio comparativo de ambas legislaciones permite identificar similitudes, diferencias y particularidades en la regulación de estas conductas delictivas.

Como parte de este ejercicio de análisis jurídico, se desarrolla un juego de roles en el que actúo como asesor jurídico, brindando orientación legal a tres consultantes que presentan situaciones hipotéticas relacionadas con estos delitos. Este ejercicio permite aplicar la teoría del delito al análisis de casos concretos, identificando elementos típicos, antijurídicos y culpables en cada supuesto planteado.


\section{Desarrollo}

\subsection{Marco normativo de los delitos contra la paz y la seguridad}

Los delitos contra la paz, la seguridad de las personas y la inviolabilidad del domicilio se encuentran regulados en diferentes ordenamientos penales de México. A continuación, se presenta un análisis comparativo de las cuatro figuras delictivas objeto de estudio:

\subsubsection{Amenazas}

\textbf{Código Penal Federal:} El artículo 282 establece: \emph{``Se aplicará sanción de tres días a un año de prisión o de 180 a 360 días multa: I.- Al que de cualquier modo amenace a otro con causarle un mal en su persona, en sus bienes, en su honor o en sus derechos, o en la persona, honor, bienes o derechos de alguien con quien esté ligado con algún vínculo, y II.- Al que por medio de amenazas de cualquier género trate de impedir que otro ejecute lo que tiene derecho a hacer.''}

\textbf{Código Penal para el Distrito Federal:} El artículo 209 señala: \emph{``Al que por cualquier medio amenace a otro con causarle un mal que, atendiendo a su contenido y contexto, sea de posible ejecución y, de causarlo, ponga en riesgo a la persona en su vida, bienes, honor o derechos, o en la persona, honor, bienes o derechos de alguien con quien esté ligado por algún vínculo, se le impondrá de uno a dos años de prisión o de noventa a trescientos sesenta días multa.''}

\subsubsection{Cobranza ilegítima}

\textbf{Código Penal Federal:} El artículo 284 Bis establece: \emph{``Se sancionará de uno a cuatro años de prisión y multa de cincuenta mil a trescientos mil pesos a quien lleve a cabo la actividad de cobranza extrajudicial ilegal. Se entiende por cobranza extrajudicial ilegal el uso de la violencia o la intimidación ilícitos, ya sea personalmente o a través de cualquier medio, para requerir el pago de una deuda derivada de actividades reguladas en leyes federales.''}

\textbf{Código Penal para el Distrito Federal:} El artículo 209 Bis (actualmente derogado, pero se conservaba la regulación en versiones anteriores) contemplaba conductas similares, aunque actualmente la Ciudad de México no tiene una figura específica denominada así, sino que aborda estas conductas a través de amenazas agravadas conforme al artículo 209 Ter.

\subsubsection{Allanamiento de morada}

\textbf{Código Penal Federal:} El artículo 285 dispone: \emph{``Se impondrán de un mes a dos años de prisión y multa de diez a cien pesos, al que, sin motivo justificado, sin orden de autoridad competente y fuera de los casos en que la ley lo permita, se introduzca, furtivamente o con engaño o violencia, o sin permiso de la persona autorizada para darlo, a un departamento, vivienda, aposento o dependencias de una casa habitada.''}

\textbf{Código Penal para el Distrito Federal:} El artículo 210 establece: \emph{``Al que se introduzca a un departamento, vivienda, aposento o dependencia de una casa habitación sin motivo justificado, sin orden de autoridad competente, furtivamente, con engaño, violencia o sin permiso de la persona autorizada para darlo, se le impondrá de seis meses a dos años de prisión o de cincuenta a cien días multa.''}

\subsubsection{Usurpación de identidad}

\textbf{Código Penal Federal:} No existe un tipo penal específico con esta denominación en el CPF. Las conductas relacionadas pueden encuadrarse en otros delitos como fraude o falsificación de documentos.

\textbf{Código Penal para el Distrito Federal:} El artículo 211 Bis señala: \emph{``Al que por cualquier medio usurpe, con fines ilícitos, la identidad de otra persona, u otorgue su consentimiento para llevar a cabo la usurpación en su identidad, se le impondrá una pena de uno a cinco años de prisión y de cuatrocientos a seiscientos días multa.''}


\subsection{Tabla comparativa de figuras delictivas}

A continuación se presenta una tabla comparativa de las cuatro figuras delictivas estudiadas, identificando la ley aplicable, artículo específico, conducta prevista y elemento principal a probar:

\begin{landscape}
\thispagestyle{empty}
\begin{longtable}{|>{\raggedright\arraybackslash}p{3.5cm}|>{\raggedright\arraybackslash}p{3cm}|>{\raggedright\arraybackslash}p{8cm}|>{\raggedright\arraybackslash}p{5cm}|}
\caption{Tabla comparativa de delitos contra la paz, seguridad e inviolabilidad del domicilio} \\
\hline
\textbf{Figura Delictiva} & \textbf{Ley / Artículo} & \textbf{Conducta Prevista} & \textbf{Elemento Principal a Probar} \\
\hline
\endfirsthead

\multicolumn{4}{c}{\tablename\ \thetable\ -- Continuación} \\
\hline
\textbf{Figura Delictiva} & \textbf{Ley / Artículo} & \textbf{Conducta Prevista} & \textbf{Elemento Principal a Probar} \\
\hline
\endhead

\hline
\multicolumn{4}{r}{\textit{Continúa en la siguiente página}} \\
\endfoot

\hline
\endlastfoot

\textbf{Amenazas} & CPF art. 282\newline CPDF art. 209 & 
Amenazar a otro por cualquier medio con causarle un mal en su persona, bienes, honor o derechos, o en personas vinculadas. El mal debe ser de posible ejecución (CPDF). Incluye amenazas para impedir que alguien ejerça un derecho. & 
La existencia y contenido de la amenaza (testigos, mensajes, grabaciones); el contexto que demuestre la posibilidad de ejecución del mal; el estado de temor o perturbación generado en la víctima; y el vínculo entre amenazador y víctima. \\
\hline

\textbf{Cobranza ilegítima / Cobranza extrajudicial ilegal} & CPF art. 284 Bis\newline CPDF art. 209 Ter (conductas relacionadas) & 
Llevar a cabo cobranza extrajudicial mediante violencia o intimidación ilícita para requerir el pago de deudas derivadas de actividades federales (créditos, financiamientos). Se agrava con uso de documentos o sellos falsos. & 
Acreditar la existencia de la deuda; demostrar actos de violencia o intimidación en el proceso de cobro (amenazas, visitas intimidatorias, contacto con terceros); testimonios de la víctima y terceros; evidencia de medios utilizados (mensajes, llamadas). \\
\hline

\textbf{Allanamiento de morada} & CPF art. 285\newline CPDF art. 210 & 
Introducirse sin motivo justificado, sin orden de autoridad competente y sin permiso de persona autorizada, en departamento, vivienda, aposento o dependencias de casa habitada, ya sea furtivamente, con engaño o violencia. & 
La introducción efectiva al domicilio ajeno; la ausencia de autorización o consentimiento del morador; la ausencia de orden judicial o causa legal justificada; testigos presenciales; evidencia de ingreso furtivo, violento o engañoso (daños, grabaciones). \\
\hline

\textbf{Usurpación de identidad} & CPDF art. 211 Bis\newline (CPF: no contempla figura específica) & 
Usurpar por cualquier medio, con fines ilícitos, la identidad de otra persona, u otorgar consentimiento para que se usurpe la propia identidad con fines ilícitos. & 
El uso indebido de datos personales identificatorios de la víctima (nombre, documentos, credenciales); la finalidad ilícita perseguida; el perjuicio o riesgo generado a la víctima; documentos, contratos o registros que evidencien la suplantación; testimonios y peritajes. \\
\hline

\end{longtable}
\end{landscape}


\subsection{Consultas jurídicas simuladas: rol de asesor jurídico}

En mi calidad de asesor jurídico, se presentan a continuación tres consultas simuladas junto con mi dictamen legal correspondiente:

\subsubsection{Consulta 1: Amenazas en contexto de conflicto vecinal}

\textbf{Consulta:} Mi vecino ha venido enviándome mensajes de texto donde me dice que ``si no dejo de estacionarme frente a su casa, va a quemar mi automóvil'' y que ``me va a partir la cara''. Tengo miedo. ¿Esto es delito?

\textbf{Dictamen jurídico:} Efectivamente, la conducta descrita constituye el delito de \textbf{amenazas} previsto en el artículo 282 del Código Penal Federal y en el artículo 209 del Código Penal para el Distrito Federal. Los mensajes de texto que menciona contienen amenazas directas a su patrimonio (quemar el automóvil) y a su integridad física (``partir la cara''), cumpliendo el tipo penal al anunciar males de posible ejecución. Recomiendo que conserve los mensajes como prueba y presente denuncia formal ante el Ministerio Público. Es importante señalar que este delito, en la mayoría de los casos, se persigue por querella, por lo que su manifestación de voluntad de proceder penalmente es indispensable. La afectación a su tranquilidad y el temor fundado que menciona son elementos que fortalecen la acreditación del delito.

\subsubsection{Consulta 2: Cobranza extrajudicial con violencia}

\textbf{Consulta:} Debo dinero a una institución de crédito. Han contratado a unos cobradores que llegan a mi casa, tocan la puerta de manera agresiva, gritan frente a mis vecinos y me han dicho que si no pago ``las cosas se pondrán feas para mi familia''. ¿Pueden hacer esto?

\textbf{Dictamen jurídico:} Lo que está experimentando configura el delito de \textbf{cobranza extrajudicial ilegal} establecido en el artículo 284 Bis del Código Penal Federal. La conducta de los cobradores constituye intimidación ilícita, al utilizar violencia moral (amenazas implícitas a su familia) y actos intimidatorios (gritos, tocamientos agresivos) para exigir el pago de una deuda. La ley es clara: si bien los acreedores tienen derecho a cobrar, no pueden utilizar violencia ni intimidación. Le sugiero documentar estos hechos (grabaciones, testigos vecinales), presentar denuncia penal ante el Ministerio Público y notificar a la Comisión Nacional para la Protección y Defensa de los Usuarios de Servicios Financieros (CONDUSEF). Estos actos constituyen delito federal sancionado con uno a cuatro años de prisión.

\subsubsection{Consulta 3: Allanamiento de morada por persona desconocida}

\textbf{Consulta:} El pasado fin de semana, mientras estaba en mi casa, una persona desconocida entró sin permiso saltando la barda del jardín. Cuando lo confronté, dijo que buscaba a alguien que no vive aquí y se retiró rápidamente. ¿Debo denunciar?

\textbf{Dictamen jurídico:} La conducta descrita encuadra en el delito de \textbf{allanamiento de morada} previsto en el artículo 285 del Código Penal Federal y artículo 210 del Código Penal para el Distrito Federal. El sujeto se introdujo a su domicilio sin autorización, de manera furtiva (saltando la barda), sin orden de autoridad competente ni causa justificada. El hecho de que alegara buscar a alguien no justifica la intromisión en su propiedad privada. Recomiendo encarecidamente presentar denuncia, ya que esta conducta viola la inviolabilidad del domicilio constitucionalmente protegida. Si cuenta con cámaras de seguridad, vecinos testigos o puede describir físicamente al intruso, estos elementos fortalecerán la carpeta de investigación. Este delito protege su derecho fundamental a la privacidad y seguridad en su hogar.


\section{Conclusiones}

Tras el análisis de los delitos contra la paz, la seguridad de las personas y la inviolabilidad del domicilio, puedo concluir que estas figuras típicas representan pilares fundamentales en la protección de la tranquilidad y seguridad individual en nuestro sistema jurídico penal. La comparativa entre el Código Penal Federal y el Código Penal para el Distrito Federal revela que, si bien ambos ordenamientos protegen bienes jurídicos similares, existen diferencias importantes en cuanto a penalidades y elementos específicos del tipo. Destaca particularmente que la usurpación de identidad, delito de creciente relevancia en la era digital, cuenta con regulación específica únicamente en el ámbito local de la Ciudad de México, evidenciando una laguna legislativa en el orden federal.

El ejercicio de juego de roles como asesor jurídico me permitió aplicar el conocimiento teórico a situaciones prácticas, identificando los elementos típicos de cada delito y los medios probatorios necesarios para su acreditación. Resulta evidente que en la práctica jurídica, el asesoramiento oportuno y preciso puede hacer la diferencia entre la protección efectiva de los derechos de las víctimas y la impunidad. Particularmente relevante es la orientación sobre la naturaleza querellable de la mayoría de estos delitos, lo que implica que sin la manifestación de voluntad de la víctima de proceder penalmente, la acción penal no puede iniciarse.

Finalmente, considero que el estudio de estos tipos penales no solo tiene valor académico, sino también un profundo sentido social. Las amenazas, la cobranza ilegítima, el allanamiento de morada y la usurpación de identidad son conductas que afectan cotidianamente a numerosas personas en nuestro país. Como futuros abogados, tenemos la responsabilidad no solo de conocer la legislación aplicable, sino de contribuir a que las personas conozcan sus derechos y los mecanismos legales disponibles para su protección. El derecho penal, en este ámbito, cumple una función esencial de tutela de la convivencia pacífica y la seguridad ciudadana.


\end{document}
